\documentclass{beamer}
\usepackage{ctex}
\usepackage[T1]{fontenc}

% other packages
\usepackage{latexsym,xcolor,multicol,booktabs,calligra}
\usepackage{pstricks,listings,stackengine,cite}
\usepackage{anyfontsize}
\usepackage{hyperref}
\usepackage{graphicx}
\usepackage{amsmath,amsthm}

\setbeamertemplate{bibliography item}[text]
% \bibliographystyle{plain}
% 如果参考文献太多的话，可以像下面这样调整字体：
% \tiny\bibliographystyle{plain}

% \usefonttheme[onlymath]{serif}

% \newtheorem{theorem}{定理}[subsection]
% \newtheorem{lemma}{引理}[subsection]
% \newtheorem*{definition}{定义}
% \newtheorem{property}{性质}[subsection]
% \newtheorem{infer}{推论}[subsection]

\usepackage{color}

\newcommand{\bsym}[1]{\boldsymbol{#1}}
\newcommand{\mypar}[1]{\left( #1 \right)}
\newcommand{\gram}[1]{\bsym{G}\mypar{#1}}
\newcommand{\abs}[1]{\left|#1 \right|}
\newcommand{\Abs}[1]{\left\Vert#1\right\Vert}
\newcommand{\setof}[1]{\left\{#1 \right\}}
\newcommand{\indot}[2]{\left\langle #1,#2 \right\rangle}
\newcommand{\mat}[1]{\left[ #1 \right]}
\newcommand{\myvec}[1]{\left[ #1 \right]^\top}
\newcommand{\sqmat}[3]{\begin{#1}
{#2}_{11} & {#2}_{12} & \cdots & {#2}_{1{#3}} \\
{#2}_{21} & {#2}_{22} & \cdots & {#2}_{2{#3}} \\
\vdots & \vdots &   \ddots     & \vdots \\
{#2}_{{#3}1} & {#2}_{{#3}2} & \cdots & {#2}_{{#3}{#3}}
\end{#1}}
\newcommand{\normmat}[4]{\begin{#1}
{#2}_{11} & {#2}_{12} & \cdots & {#2}_{1{#4}} \\
{#2}_{21} & {#2}_{22} & \cdots & {#2}_{2{#4}} \\
\vdots & \vdots &   \ddots     & \vdots \\
{#2}_{{#3}1} & {#2}_{{#3}2} & \cdots & {#2}_{{#3}{#4}}
\end{#1}}
\newcommand{\nullsp}[1]{\bsym{\mathrm{N}}\mypar{#1}}
\newcommand{\colsp}[1]{\bsym{\mathrm{C}}\mypar{#1}}
\newcommand{\func}[2]{\mathrm{#1}\mypar{#2}}
\newcommand{\entry}[3]{\func{entry}{#1,#2,#3}}
\newcommand{\row}[2]{\func{row}{#1,#2}}
\newcommand{\col}[2]{\func{col}{#1,#2}}
\newcommand{\trace}[1]{\func{Tr}{#1}}
\newcommand{\diag}[1]{\func{diag}{#1}}
\newcommand{\rank}[1]{\func{rank}{#1}}
\newcommand{\adj}[1]{\func{adj}{#1}}
\newcommand{\myspan}[1]{\func{span}{#1}}
\newcommand{\enums}[2]{{#1}_1,{#1}_2,\dots,{#1}_{#2}}
\newcommand{\inv}[1]{{#1}^{-1}}
\newcommand{\pinv}[1]{\inv{\mypar{#1}}}
\newcommand{\ortcom}[1]{{#1}^{\bot}}

\renewcommand{\det}[1]{\func{det}{#1}}

\newcommand{\dif}{\mathrm{d}}
\newcommand{\fracdif}[1]{\frac{\dif}{\dif #1}}

\newcommand{\todo}[1]{{ \textcolor{red}{ TODO: #1}}}

\newcommand{\mata}{\bsym{A}}
\newcommand{\matb}{\bsym{B}}
\newcommand{\matc}{\bsym{C}}
\newcommand{\matd}{\bsym{D}}
\newcommand{\mate}{\bsym{E}}
\newcommand{\matf}{\bsym{F}}
\newcommand{\matfstar}{\bsym{F^*}}
\newcommand{\matg}{\bsym{G}}
\newcommand{\matH}{\bsym{H}}
\newcommand{\mati}{\bsym{I}}
\newcommand{\matj}{\bsym{J}}
\newcommand{\matk}{\bsym{K}}
\newcommand{\matl}{\bsym{L}}
\newcommand{\matlam}{\bsym{\Lambda}}
\newcommand{\matm}{\bsym{M}}
\newcommand{\matn}{\bsym{N}}
\newcommand{\mato}{\bsym{O}}
\newcommand{\matp}{\bsym{P}}
\newcommand{\matq}{\bsym{Q}}
\newcommand{\matr}{\bsym{R}}
\newcommand{\mats}{\bsym{S}}
\newcommand{\matsig}{\bsym{\Sigma}}
\newcommand{\matt}{\bsym{T}}
\newcommand{\matu}{\bsym{U}}
\newcommand{\matv}{\bsym{V}}
\newcommand{\matw}{\bsym{W}}
\newcommand{\matx}{\bsym{X}}
\newcommand{\maty}{\bsym{Y}}
\newcommand{\matz}{\bsym{Z}}

\newcommand{\field}{\bsym{\mathrm{F}}}
\newcommand{\rea}{\bsym{\mathbb{R}}}

\newcommand{\veca}{\bsym{a}}
\newcommand{\vecal}{\bsym{\alpha}}
\newcommand{\vecb}{\bsym{b}}
\newcommand{\vecbeta}{\bsym{\beta}}
\newcommand{\vecc}{\bsym{c}}
\newcommand{\vecd}{\bsym{d}}
\newcommand{\vecdelta}{\bsym{\delta}}
\newcommand{\vece}{\bsym{e}}
\newcommand{\veceps}{\bsym{\epsilon}}
\newcommand{\vecveps}{\bsym{\varepsilon}}
\newcommand{\vecf}{\bsym{f}}
\newcommand{\vecg}{\bsym{g}}
\newcommand{\vecgamma}{\bsym{\gamma}}
\newcommand{\vech}{\bsym{h}}
\newcommand{\veceta}{\bsym{\eta}}
\newcommand{\veci}{\bsym{i}}
\newcommand{\vecj}{\bsym{j}}
\newcommand{\veck}{\bsym{k}}
\newcommand{\vecl}{\bsym{l}}
\newcommand{\vecm}{\bsym{m}}
\newcommand{\vecn}{\bsym{n}}
\newcommand{\veco}{\bsym{o}}
\newcommand{\vecone}{\bsym{1}}
\newcommand{\vecp}{\bsym{p}}
\newcommand{\vecq}{\bsym{q}}
\newcommand{\vecr}{\bsym{r}}
\newcommand{\vecs}{\bsym{s}}
\newcommand{\vect}{\bsym{t}}
\newcommand{\vecu}{\bsym{u}}
\newcommand{\vecv}{\bsym{v}}
\newcommand{\vecw}{\bsym{w}}
\newcommand{\vecx}{\bsym{x}}
\newcommand{\vecxi}{\bsym{\xi}}
\newcommand{\vecy}{\bsym{y}}
\newcommand{\vecz}{\bsym{z}}
\newcommand{\veczero}{\bsym{0}}


\author{夏海淞}
\title{线性代数习题课}
\subtitle{第五章}
% \institute{School of Computer Science, Fudan University}
\date{\today}
\usepackage{College}

% defs
\def\cmd#1{\texttt{\color{red}\footnotesize $\backslash$#1}}
\def\env#1{\texttt{\color{blue}\footnotesize #1}}
\definecolor{deepblue}{rgb}{0,0,0.5}
\definecolor{deepred}{rgb}{0.6,0,0}
\definecolor{deepgreen}{rgb}{0,0.5,0}
\definecolor{halfgray}{gray}{0.55}

\lstset{
    basicstyle=\ttfamily\small,
    keywordstyle=\bfseries\color{deepblue},
    emphstyle=\ttfamily\color{deepred},    % Custom highlighting style
    stringstyle=\color{deepgreen},
    numbers=left,
    numberstyle=\small\color{halfgray},
    rulesepcolor=\color{red!20!green!20!blue!20},
    frame=shadowbox,
}

\begin{document}
\setlength{\parskip}{0.45\baselineskip}

\begin{frame}
    \titlepage
    \begin{figure}[htpb]
        \begin{center}
            \vspace*{-0.5cm}
            \includegraphics[width=0.1\linewidth]{../../pic/FDU.jpeg}
        \end{center}
    \end{figure}
\end{frame}

\subsection*{5.1}
\begin{frame}
    \frametitle{题面}
    求下列矩阵的特征值和特征向量：
    \begin{enumerate}
        \item[(2)]{
              \begin{equation*}
                  \begin{bmatrix}
                      -1 & 2  & 2  \\
                      3  & -1 & 1  \\
                      2  & 2  & -1
                  \end{bmatrix}
              \end{equation*}
              }
    \end{enumerate}
\end{frame}

\begin{frame}
    \frametitle{解答}
    记上述矩阵为\(\mata\)，则矩阵\(\mata\)的特征多项式为
    \begin{equation*}
        \abs{\lambda\mati-\mata} =
        \begin{vmatrix}
            \lambda+1 & -2        & -2        \\
            -3        & \lambda+1 & -1        \\
            -2        & -2        & \lambda+1
        \end{vmatrix}=\spar{\lambda+3}^2\spar{\lambda-3}
    \end{equation*}
    因此\(\mata\)的特征值为\(\lambda_1=\lambda_2=-3,\lambda_3=3\)。
    \pause

    将\(\lambda_1=\lambda_2=-3\)代入\(\spar{\lambda\mati-\mata}\vecx=\veczero\)，得到其基础解系为\(\myvec{1,-2,1}\)。因此属于特征值\(-3\)的特征向量为\(k_1\myvec{1,-2,1},k_1\neq0\)。

    将\(\lambda_3=3\)代入\(\spar{\lambda\mati-\mata}\vecx=\veczero\)，得到其基础解系为\(\myvec{1,1,1}\)。因此属于特征值\(3\)的特征向量为\(k_2\myvec{1,1,1},k_2\neq0\)。
\end{frame}

\subsection*{5.2}
\begin{frame}
    \frametitle{题面}
    设矩阵\(\mata=\begin{bmatrix}a&0&b\\0&2&0\\b&0&-2\end{bmatrix}\)的一个特征值为\(\lambda_1=-3\)，且三个特征值的积为\(-12\)，则求\(a\)，\(b\)，以及\(\mata\)的其他特征值。
\end{frame}

\begin{frame}
    \frametitle{解答}
    由特征向量的定义和定理5.1.2，有
    \begin{equation*}
        \begin{cases}
            \lambda_1+\lambda_2+\lambda_3=\trace{\mata}=a        \\
            \lambda_1\lambda_2\lambda_3=\det{\mata}=-4a-2b^2=-12 \\
            \det{\lambda_1\mati-\mata}=-5a+5b^2-15=0
        \end{cases}
    \end{equation*}
    \pause
    解得\(a=1\)，\(b=\pm2\)，\(\lambda_2=\lambda_3=2\)。
\end{frame}

\subsection*{5.3}
\begin{frame}
    \frametitle{题面}
    已知\(\vecal=\myvec{1,1,-1}\)是矩阵
    \begin{equation*}
        \mata=
        \begin{bmatrix}
            2  & -1 & 2  \\
            5  & a  & 3  \\
            -1 & b  & -2
        \end{bmatrix}
    \end{equation*}
    的一个特征向量，试确定\(a\)，\(b\)的值以及特征向量\(\vecal\)所对应的特征值。
\end{frame}

\begin{frame}
    \frametitle{解答}
    设\(\vecal\)对应的特征值为\(\lambda\)，则由特征值和特征向量的定义可知\(\mata\vecal=\lambda\vecal\)。将\(\mata,\vecal\)代入即可解得\(\lambda=-1\)，\(a=-3\)，\(b=0\)。
\end{frame}

\subsection*{5.4}
\begin{frame}
    \frametitle{题面}
    设\(3\)阶矩阵\(\mata\)的三个特征值为\(\lambda_1=1\)，\(\lambda_2=2\)，\(\lambda_3=3\)，与之对应的特征向量分别为\(\vecal_1=\myvec{2,1,-1}\)，\(\vecal_2=\myvec{2,-1,2}\)，\(\vecal_=\myvec{3,0,1}\)，求矩阵\(\mata\)。
\end{frame}

\begin{frame}
    \frametitle{解答}
    由特征值和特征向量的定义可知
    \begin{equation*}
        \mata\begin{bmatrix}\vecal_1&\vecal_2&\vecal_3\end{bmatrix}=\begin{bmatrix}\lambda_1\vecal_1&\lambda_2\vecal_2&\lambda_3\vecal_3\end{bmatrix}
    \end{equation*}
    \pause
    由定理5.1.5知\(\mat{\vecal_1,\vecal_2,\vecal_3}\)可逆，因此有
    \begin{equation*}
        \mata=\begin{bmatrix}\lambda_1\vecal_1&\lambda_2\vecal_2&\lambda_3\vecal_3\end{bmatrix}\inv{\begin{bmatrix}\vecal_1&\vecal_2&\vecal_3\end{bmatrix}}=
        \begin{bmatrix}
            -3 & 26 & 18 \\
            -1 & 6  & 3  \\
            0  & 2  & 3
        \end{bmatrix}
    \end{equation*}
\end{frame}

\subsection*{5.5}
\begin{frame}
    \frametitle{题面}
    设\(\mata^2=\mata\)，证明\(\mata\)的特征值只能是\(0\)或者\(1\)。
\end{frame}

\begin{frame}
    \frametitle{解答}
    设\(\mata\)的特征向量为\(\vecal\neq\veczero\)。由题设及特征向量的定义，有
    \begin{equation*}
        \lambda\vecal=\mata\vecal=\mata^2\vecal=\lambda^2\vecal
    \end{equation*}
    \pause
    因此可得\(\spar{\lambda^2-\lambda}\vecal=\veczero\)。由\(\vecal\neq\veczero\)解得\(\lambda=0\)或\(\lambda=1\)。

    结论成立。
\end{frame}

\subsection*{5.6}
\begin{frame}
    \frametitle{题面}
    记\(\mata\)的伴随矩阵为\(\mata^*\)。设\(n\)阶可逆矩阵\(\mata\)的特征值为\(n\)个非零数\(\enums{\lambda}{n}\)，试证\(\mata\)的伴随矩阵的特征值为\(\det{\mata}\inv{\lambda_1}\)，\(\det{\mata}\inv{\lambda_2}\)，\(\dots\)，\(\det{\mata}\inv{\lambda_n}\)。
\end{frame}

\begin{frame}
    \frametitle{解答}
    对\(\forall i\in\setof{1,2,\dots,n}\)，设\(\vecal_i\)为\(\lambda_i\)对应的一个特征向量。则由定理3.3.1和特征向量的定义，有
    \begin{equation*}
        \det{\mata}\vecal=\mata^*\mata\vecal=\lambda_i\mata^*\vecal
    \end{equation*}
    \pause
    所以\(\det{\mata}\inv{\lambda_i}\spar{i=1,2,\dots,n}\)是\(\mata^*\)的特征值。

    由\(\mata^*\in\rea^{n\times n}\)可知\(\mata^*\)没有其他特征值。
\end{frame}

\subsection*{5.7}
\begin{frame}
    \frametitle{题面}
    已知\(3\)阶方阵\(\mata\)的特征值为\(1,2,-1\)，设矩阵\(\matb=\mata^3-5\mata^2\)。试计算行列式\(\det{\matb}\)，\(\det{\mata-5\mati}\)。
\end{frame}

\begin{frame}
    \frametitle{解答}

    由推论5.2.1可知存在可逆矩阵\(\matp\)，满足\(\mata=\matp\diag{\lambda_1,\lambda_2,\lambda_3}\inv{\matp}\)。因此有
    \begin{align*}
        \mata-5\mati & =\matp\diag{\lambda_1,\lambda_2,\lambda_3}\inv{\matp}-5\matp\mati\inv{\matp} \\
                     & =\matp\diag{\lambda_1-5,\lambda_2-5,\lambda_3-5}\inv{\matp}
    \end{align*}
    \pause
    所以\(\det{\mata-5\mati}=\prod_{i=1}^3\spar{\lambda_i-5}=-72\)。

    同理可得\(\det{\matb}=\prod_{i=1}^3\spar{\lambda_i^3-5\lambda_i^2}=-288\)。

\end{frame}

\subsection*{5.8}
\begin{frame}
    \frametitle{题面}
    设\(\mata=\mati-2\vecx\vecx^\top\)，其中\(\vecx=\myvec{\enums{x}{n}}\)，\(\mati\)为\(n\)阶单位矩阵。若\(\vecx^\top\vecx=1\)，求证：\(\vecx\)为\(\mata\)的属于特征值\(-1\)的特征向量。
\end{frame}

\begin{frame}
    \frametitle{解答}
    由\(\vecx^\top\vecx=1\)可得
    \begin{align*}
        \mata\vecx & =\spar{\mati-2\vecx\vecx^\top}\vecx \\
                   & =\mati\vecx-2\vecx\vecx^\top\vecx   \\
                   & =\vecx-2\vecx=-\vecx
    \end{align*}
    \pause
    由特征值和特征向量的定义可知\(\vecx\)为\(\mata\)的属于特征值\(-1\)的特征向量。
\end{frame}

\subsection*{5.9}
\begin{frame}
    \frametitle{题面}
    已知\(n\)阶矩阵\(\mata\)的特征值为\(\lambda_0\)，
    \begin{enumerate}
        \item 求\(k\mata\)的特征值（\(k\)为任意实数）；
        \item 若\(\mata\)可逆，求\(\inv{\mata}\)的特征值；
        \item 求\(\mati+\mata\)的特征值。
    \end{enumerate}
\end{frame}

\begin{frame}
    \frametitle{解答}
    设\(\mata\)的特征值为\(\lambda_0\)对应的特征向量为\(\vecal_0\)。
    \begin{enumerate}
        \item 当\(k\neq0\)时，有\(\mata\vecal_0=\lambda_0\vecal_0\iff k\mata\vecal_0=k\lambda_0\vecal_0\)。此时\(k\mata\)的特征值为\(k\lambda_0\)。当\(k=0\)时结论依然成立。\pause
        \item 因为\(\mata\)可逆，所以\(\lambda_0\neq0\)。因此有\(\mata\vecal_0=\lambda_0\vecal_0\iff\inv{\lambda_0}\vecal_0=\inv{\mata}\vecal_0\)。所以\(\inv{\mata}\)的特征值为\(\inv{\lambda_0}\)。\pause
        \item 同理，可得\(\mata\vecal_0=\lambda_0\vecal_0\iff\spar{\mata+\mati}\vecal_0=\spar{\lambda_0+1}\vecal_0\)。所以\(\mati+\mata\)的特征值为\(\lambda_0+1\)。
    \end{enumerate}
\end{frame}

\subsection*{5.10}
\begin{frame}
    \frametitle{题面}
    设\(\mata^2-3\mata+2\mati=\mato\)，证明\(\mata\)的特征值只能是\(1\)或\(2\)。
\end{frame}

\begin{frame}
    \frametitle{解答}
    设\(\mata\)的特征向量为\(\vecal\neq\veczero\)。由题设及特征向量的定义，有
    \begin{equation*}
        \spar{\mata^2-3\mata+2\mati}\vecal=\spar{\lambda^2-3\lambda+2}\vecal=\veczero
    \end{equation*}
    \pause
    由\(\vecal\neq\veczero\)解得\(\lambda=1\)或\(\lambda=2\)。结论成立。
\end{frame}

\subsection*{5.12}
\begin{frame}
    \frametitle{题面}
    设非零向量\(\vecal=\myvec{\enums{a}{n}}\)，\(\vecbeta=\myvec{\enums{b}{n}}\)，且\(\vecal^\top\vecbeta=0\)，\(\mata=\vecal\vecbeta^\top\)，求\(\mata\)的特征值和特征向量。
\end{frame}

\begin{frame}[allowframebreaks]
    \frametitle{解答}

    设\(\mata\)的特征值为\(\lambda\)，对应的特征向量为\(\vecx\neq\veczero\)。由\(\vecal^\top\vecbeta=0\)和特征值的定义可得
    \begin{equation*}
        \lambda^2\vecx=\mata^2\vecx=\vecal\spar{\vecbeta^\top\vecal}\vecbeta^\top\vecx=\veczero
    \end{equation*}
    解得\(\lambda=0\)，即\(\mata\)的特征值均为\(0\)。

    又因为\(\vecal,\vecbeta\)为非零向量，所以\(\rank{\mata}=\rank{\vecal\vecbeta^\top}=\rank{\vecbeta^\top}\)。因此方程\(\vecal\vecbeta^\top\vecx\)和方程\(\vecbeta^\top\vecx\)同解。所以\(\mata\)的特征向量为方程
    \begin{equation*}
        \vecbeta^\top\vecx=b_1x_1+b_2x_2+\dots+b_nx_n=0
    \end{equation*}
    的解。不妨设\(b_i\neq0\)，则\(\mata\)线性无关的特征向量分别为
    \begin{gather*}
        \begin{bmatrix}b_i&0&\cdots&0&-b_1&0&\cdots&0&0\end{bmatrix}^\top\\
        \begin{bmatrix}0&b_i&\cdots&0&-b_2&0&\cdots&0&0\end{bmatrix}^\top\\
        \vdots\\
        \begin{bmatrix}0&0&\cdots&0&-b_n&0&\cdots&0&b_i\end{bmatrix}^\top
    \end{gather*}

\end{frame}

\subsection*{5.13}
\begin{frame}
    \frametitle{题面}
    已知\(\lambda_1=2\)是\(3\)阶矩阵\(\mata\)的一个特征值，\(\vecal_1=\mat{1,2,0}\)，\(\vecal_2=\mat{1,0,1}\)是\(\mata\)属于特征值\(2\)的两个特征向量，\(\vecbeta=\myvec{-1,2,-2}\)，计算\(\mata\vecbeta\)。
\end{frame}

\begin{frame}
    \frametitle{解答}
    容易发现\(\vecal_1,\vecal_2,\vecbeta\)线性相关：\(\vecbeta=\vecal_1-2\vecal_2\)。因此由特征值和特征向量的定义，有
    \begin{equation*}
        \mata\vecbeta=\mata\spar{\vecal_1-2\vecal_2}=2\spar{\vecal_1-2\vecal_2}=\begin{bmatrix}2&4&-4\end{bmatrix}^\top
    \end{equation*}
\end{frame}

\subsection*{5.14}
\begin{frame}
    \frametitle{题面}
    设矩阵\(\mata=\begin{bmatrix}a&-1&c\\5&b&3\\1-c&0&-a\end{bmatrix}\)，且\(\det{\mata}=-1\)。伴随矩阵\(\mata^*\)有特征值\(\lambda_0\)，属于\(\lambda_0\)的特征向量为\(\vecal=\myvec{-1,-1,1}\)，求\(a,b,c\)，以及\(\lambda_0\)。
\end{frame}

\begin{frame}
    \frametitle{解答}
    由题意可得\(\mata^*\vecal=\lambda_0\vecal\)，通过定理3.3.1可转化为
    \begin{equation*}
        \mata\vecal=\inv{\lambda_0}\det{\mata}\vecx=-\inv{\lambda_0}\vecx
    \end{equation*}
    \pause
    将其与\(\det{\mata}=-1\)联立可得
    \begin{equation*}
        \begin{cases}
            -a^2b-\spar{3+bc}\spar{1-c}-5a=-1 \\
            -a+1+c=-\inv{\lambda_0}           \\
            -b-2=-\inv{\lambda_0}             \\
            c-1-a=-\inv{\lambda_0}
        \end{cases}
    \end{equation*}
    \pause
    解得\(a=c=2,b=-3,\lambda_0=1\)。
\end{frame}

\subsection*{5.15}
\begin{frame}
    \frametitle{题面}
    若\(n\)阶方阵\(\mata\)的每行元素之和为常数\(c\)，证明：\(c\)是\(\mata\)的一个特征值：\(\forall m\in\mathbf{N}\)，\(\mata^m\)的每行元素之和为\(c^m\)。
\end{frame}

\begin{frame}
    \frametitle{解答}
    记\(\rea^n\)中各分量均为\(1\)的列向量为\(\vecone\)。则由题设可知\(\mata\vecone=c\vecone\)。
    \pause

    所以\(c\)是\(\mata\)的一个特征值，\(\vecone\)是其对应的特征向量。
    \pause

    因此有\(\mata^m\vecone=c^m\vecone\)，结论成立。
\end{frame}

\subsection*{5.18}
\begin{frame}
    \frametitle{题面}
    设\(\enums{\vecx}{n}\)是\(\mata\)的属于\(\lambda_0\)的特征向量，试证\(\enums{\vecx}{n}\)的任一非零线性组合也是\(\mata\)的属于\(\lambda_0\)的特征向量。
\end{frame}

\begin{frame}
    \frametitle{解答}
    设\(\rea^n\)中的任意非零向量为\(\vecy\)，则\(\enums{\vecx}{n}\)的任一非零线性组合可表示为\(\mat{\enums{\vecx}{n}}\vecy\)。
    \pause
    因此有
    \begin{align*}
          & \mata\begin{bmatrix}\vecx_1&\vecx_2&\cdots&\vecx_n\end{bmatrix}\vecy      \\
        = & \begin{bmatrix}\mata\vecx_1&\mata\vecx_2&\cdots&\mata\vecx_n\end{bmatrix} \\
        = & \lambda_0\begin{bmatrix}\vecx_1&\vecx_2&\cdots&\vecx_n\end{bmatrix}\vecy
    \end{align*}
    结论成立。
\end{frame}

\subsection*{5.19}
\begin{frame}
    \frametitle{题面}
    设\(\vecx_1,\vecx_2\)分别是\(\mata\)的属于\(\lambda_1,\lambda_2\)的特征向量，而且\(\lambda_1\neq\lambda_2\)，试证\(\vecx_1+\vecx_2\)不可能是\(\mata\)的特征向量。
\end{frame}

\begin{frame}
    \frametitle{解答}
    反设结论不成立，即\(\vecx_1+\vecx_2\)是\(\mata\)的特征值为\(\lambda\)的特征向量。则有
    \begin{equation*}
        \lambda_1\vecx_1+\lambda_2\vecx_2=\mata\vecx_1+\mata\vecx_2=\mata\spar{\vecx_1+\vecx_2}=\lambda\spar{\vecx_1+\vecx_2}
    \end{equation*}
    \pause
    化简得到\(\spar{\lambda_1-\lambda}\vecx_1+\spar{\lambda_2-\lambda}\vecx_2=\veczero\)。由\(\lambda_1\neq\lambda_2\)可知\(\vecx_1,\vecx_2\)线性相关，和定理5.1.5矛盾。

    所以\(\vecx_1+\vecx_2\)不可能是\(\mata\)的特征向量。
\end{frame}

\subsection*{5.24}
\begin{frame}
    \frametitle{题面}
    设\(\mata=\vecu\vecv^\top\)，其中\(\vecu,\vecv\)都是\(n\)元非零列向量，求证：\(\vecu\)是\(\mata\)的特征向量，并求出其对应的特征值。
\end{frame}

\begin{frame}
    \frametitle{解答}
    因为\(\mata\vecu=\vecu\vecv^\top\vecu=\spar{\vecu^\top\vecv}\vecu\)，由特征值和特征向量的定义可知\(\vecu\)是\(\mata\)的特征向量，对应的特征值是\(\vecu^\top\vecv\)。
\end{frame}

\subsection*{5.25}
\begin{frame}
    \frametitle{题面}
    设\(\mata=\begin{bmatrix}2&0&0\\0&0&1\\0&1&x\end{bmatrix}\)，\(\matb=\begin{bmatrix}2&&\\&y&\\&&-1\end{bmatrix}\)，已知\(\mata\)与\(\matb\)相似，求\(x,y\)。
\end{frame}

\begin{frame}
    \frametitle{解答}
    由定理5.1.2和定理5.1.3可知，相似矩阵的迹和行列式相等。因此有
    \begin{equation*}
        \begin{cases}
            -2=-2y \\
            2+x=1+y
        \end{cases}
    \end{equation*}
    解得\(x=0,y=1\)。
\end{frame}

\subsection*{5.26}
\begin{frame}
    \frametitle{题面}
    设\(\mata=\begin{bmatrix}-3&2\\-2&2\end{bmatrix}\)。
    \begin{enumerate}
        \item 求可逆矩阵\(\matp\)使得\(\inv{\matp}\mata\matp\)为对角矩阵；
        \item 计算\(\mata^{10}-\mata^6-\mati\)。
    \end{enumerate}
\end{frame}

\begin{frame}[allowframebreaks]
    \frametitle{解答}
    \begin{enumerate}
        \item {
              \(\mata\)的特征多项式为
              \begin{equation*}
                  \abs{\lambda\mati-\mata}=\begin{vmatrix}\lambda+3&-2\\2&\lambda-2\end{vmatrix}=\spar{\lambda-1}\spar{\lambda+2}
              \end{equation*}
              因此\(\mata\)的特征值为\(\lambda_1=1,\lambda_2=-2\)。

              将\(\lambda_1=1\)代入特征方程组，解得其线性无关的特征向量为\(\myvec{1,2}\)；

              将\(\lambda_2=-2\)代入特征方程组，解得其线性无关的特征向量为\(\myvec{2,1}\)。

              因此得到满足条件的\(\matp=\begin{bmatrix}1&2\\2&1\end{bmatrix}\)。
              }
        \item {
              由(1)可得\(\mata=\matp\matlam\inv{\matp}\)，其中\(\matlam=\begin{bmatrix}1&0\\0&-2\end{bmatrix}\)。因此有
              \begin{align*}
                  \mata^{10}-\mata^6-\mati & =\matp\spar{\matlam^{10}-\matlam^6-\mati}\inv{\matp}                                                                       \\
                                           & =\begin{bmatrix}1&2\\2&1\end{bmatrix}\begin{bmatrix}-1&0\\0&959\end{bmatrix}\begin{bmatrix}-1/3&2/3\\2/3&-1/3\end{bmatrix} \\
                                           & =\begin{bmatrix}1279&-640\\640&-321\end{bmatrix}
              \end{align*}
              }
    \end{enumerate}
\end{frame}

\subsection*{5.27}
\begin{frame}
    \frametitle{题面}
    已知方阵\(\mata=\begin{bmatrix}1&-1&1\\x&4&y\\-3&-3&5\end{bmatrix}\)与对角矩阵相似，且\(\lambda=2\)是\(\mata\)的二重特征值。
    \begin{enumerate}
        \item 求\(x\)与\(y\)的值；
        \item 求可逆矩阵\(\matp\)使得\(\inv{\matp}\mata\matp\)为对角阵。
    \end{enumerate}
\end{frame}

\begin{frame}
    \frametitle{解答}
    \begin{enumerate}
        \item {
              因为\(\mata\)可对角化，由定理5.2.2可知\(\mata\)的代数重数与几何重数相等。考虑
              \begin{equation*}
                  \lambda\mati-\mata=2\mati-\mata=
                  \begin{bmatrix}
                      1  & 1  & -1 \\
                      -x & -2 & -y \\
                      3  & 3  & -3
                  \end{bmatrix}\to
                  \begin{bmatrix}
                      1   & 1 & -1  \\
                      2-x & 0 & 2-y \\
                      0   & 0 & 0
                  \end{bmatrix}
              \end{equation*}
              因为\(\lambda=2\)是\(\mata\)的二重特征值，所以\(\rank{\lambda\mati-\mata}=1\)，由此得到\(x=2,y=-2\)。
              }\pause
        \item {
              由定理5.1.2可知\(\mata\)的另一个特征值为\(\trace{\mata}-2\lambda=6\)。
              将\(\lambda_1=\lambda_2=2\)，\(\lambda_3=6\)代入特征方程组\(\spar{\lambda\mati-\mata}\vecal=0\)，得到\(\mata\)对应的线性无关的特征向量分别为\(\myvec{-1,1,0}\)，\(\myvec{1,0,1}\)和\(\myvec{1,-2,3}\)。所以对应的可逆矩阵\(\matp\)为\(\begin{bmatrix}-1&1&1\\1&0&-2\\0&1&3\end{bmatrix}\)。
              }
    \end{enumerate}
\end{frame}

\subsection*{5.28}
\begin{frame}
    \frametitle{题面}
    已知矩阵\(\mata\)，求使\(\inv{\matq}\mata\matq\)为对角阵的正交矩阵\(\matq\)。
    \begin{enumerate}
        \item[(2)]
              {
              \begin{equation*}
                  \mata=
                  \begin{bmatrix}
                      3 & 2 & 4 \\
                      2 & 0 & 2 \\
                      4 & 2 & 3
                  \end{bmatrix}
              \end{equation*}
              }
    \end{enumerate}
\end{frame}

\begin{frame}[allowframebreaks]
    \frametitle{解答}
    \(\mata\)的特征多项式为
    \begin{equation*}
        \abs{\lambda\mati-\mata}=\begin{vmatrix}\lambda-1&-1&-1\\-1&\lambda-1&-1\\-1&\lambda-1&-1\end{vmatrix}=\lambda^2\spar{\lambda-3}
    \end{equation*}
    所以\(\mata\)的特征值为\(\lambda_1=\lambda_2=0\)，\(\lambda_3=3\)。

    将上述特征值代入特征方程组\(\spar{\lambda\mati-\mata}\vecal=0\)，得到\(\lambda_1,\lambda_2\)对应的线性无关的特征向量为\(\vecal_1=\myvec{-1,1,0}\)和\(\vecal_2=\myvec{-1,0,1}\)，\(\lambda_3\)对应的线性无关的特征向量为\(\vecal_3=\myvec{1,1,1}\)。

    对向量组\(\vecal_1,\vecal_2\)和向量组\(\vecal_3\)分别做Schmidt正交化，得到正交矩阵
    \begin{equation*}
        \matq=\begin{bmatrix}\vecq_1&\vecq_2&\vecq_3\end{bmatrix}=
        \begin{bmatrix}
            -\frac{\sqrt{2}}{2} & -\frac{\sqrt{6}}{6} & \frac{\sqrt{3}}{3} \\
            \frac{\sqrt{2}}{2}  & -\frac{\sqrt{6}}{6} & \frac{\sqrt{3}}{3} \\
            0                   & \frac{2\sqrt{6}}{6} & \frac{\sqrt{3}}{3}
        \end{bmatrix}
    \end{equation*}
\end{frame}

\subsection*{5.30}
\begin{frame}
    \frametitle{题面}
    设矩阵\(\mata=\begin{bmatrix}0&0&1\\a&1&b\\1&0&0\end{bmatrix}\)有三个线性无关的特征向量，求\(a,b\)满足的条件，并求可逆矩阵\(\matp\)，使得\(\inv{\matp}\mata\matp\)为对角阵。
\end{frame}

\begin{frame}[allowframebreaks]
    \frametitle{解答}
    \(\mata\)的特征多项式为
    \begin{equation*}
        \abs{\lambda\mati-\mata}=\begin{vmatrix}\lambda&0&-1\\-a&\lambda-1&-b\\-1&0&\lambda\end{vmatrix}=\spar{\lambda-1}^2\spar{\lambda+1}
    \end{equation*}
    所以\(\mata\)的特征值为\(\lambda_1=\lambda_2=1\)，\(\lambda_3=-1\)。

    由定理5.2.1和定理5.2.2可知\(\mata\)的代数重数与几何重数相等。考虑
    \begin{equation*}
        \lambda_1\mati-\mata=\mati-\mata=
        \begin{bmatrix}
            1  & 0 & -1 \\
            -a & 0 & -b \\
            -1 & 0 & 1
        \end{bmatrix}\to
        \begin{bmatrix}
            1  & 0 & -1 \\
            -a & 0 & -b \\
            0  & 0 & 0
        \end{bmatrix}
    \end{equation*}
    因为\(\lambda=1\)是\(\mata\)的二重特征值，所以\(\rank{\lambda\mati-\mata}=1\)，由此得到\(a=b=0\)。

    将\(\lambda_1=\lambda_2=1\)，\(\lambda_3=-1\)代入特征方程组\(\spar{\lambda\mati-\mata}\vecal=0\)，得到\(\mata\)对应的线性无关的特征向量分别为\(\myvec{0,1,0}\)，\(\myvec{1,0,1}\)和\(\myvec{-1,0,1}\)。所以对应的可逆矩阵\(\matp\)为\(\begin{bmatrix}0&1&-1\\1&0&0\\0&1&1\end{bmatrix}\)。
\end{frame}

\subsection*{5.32}
\begin{frame}
    \frametitle{题面}
    计算\({\begin{bmatrix}1&2&2\\2&1&2\\2&2&1\end{bmatrix}}^{100}\)。
\end{frame}

\begin{frame}[allowframebreaks]
    \frametitle{解答}
    记矩阵\(\begin{bmatrix}1&2&2\\2&1&2\\2&2&1\end{bmatrix}\)为\(\mata\)。则\(\mata\)的特征多项式为
    \begin{equation*}
        \abs{\lambda\mati-\mata}=\begin{vmatrix}\lambda-1&-2&-2\\-2&\lambda-1&-2\\-2&-2&\lambda-1\end{vmatrix}=\spar{\lambda+1}^2\spar{\lambda-5}
    \end{equation*}
    所以\(\mata\)的特征值为\(\lambda_1=\lambda_2=-1\)，\(\lambda_3=5\)。

    将\(\lambda_1=\lambda_2=-1\)，\(\lambda_3=5\)代入特征方程组\(\spar{\lambda\mati-\mata}\vecal=0\)，得到\(\mata\)对应的线性无关的特征向量分别为\(\myvec{-1,1,0}\)，\(\myvec{-1,0,1}\)和\(\myvec{1,1,1}\)。所以有
    \begin{equation*}
        \matp=\begin{bmatrix}-1&-1&1\\1&0&1\\0&1&1\end{bmatrix},\matlam=\begin{bmatrix}-1&0&0\\0&-1&0\\0&0&5\end{bmatrix}
    \end{equation*}
    \begin{align*}
        \mata^{100} & =\spar{\matp\matlam\inv{\matp}}^{100}=\matp\matlam^{100}\inv{\matp}                                                                                                   \\
                    & =\frac{1}{3}\begin{bmatrix}-1&-1&1\\1&0&1\\0&1&1\end{bmatrix}\begin{bmatrix}1&0&0\\0&1&0\\0&0&5^{100}\end{bmatrix}\begin{bmatrix}-1&2&-1\\-1&-1&2\\1&1&1\end{bmatrix} \\
                    & =\frac{1}{3}\begin{bmatrix}5^{100}+2&5^{100}-1&5^{100}-1\\5^{100}-1&5^{100}+2&5^{100}-1\\5^{100}-1&5^{100}-1&5^{100}+2\end{bmatrix}
    \end{align*}
\end{frame}

\subsection*{5.34}
\begin{frame}
    \frametitle{题面}
    求二次型\(f\)的矩阵：
    \begin{enumerate}
        \item[(2)] \(f(x_1,x_2,x_3)=\spar{x_1+x_2}^2+\spar{x_2-x_3}^2+\spar{x_3+x_1}^2\)
    \end{enumerate}
\end{frame}

\begin{frame}
    \frametitle{解答}
    \begin{align*}
        f(x_1,x_2,x_3) & =\spar{x_1+x_2}^2+\spar{x_2-x_3}^2+\spar{x_3+x_1}^2 \\
                       & =2x_1^2+2x_2^2+2x_3^2+2x_1x_2-2x_2x_3+2x_3x_1
    \end{align*}

    所以二次型\(f\)的矩阵为\(\mata=\begin{bmatrix}2&1&1\\1&2&-1\\1&-1&2\end{bmatrix}\)。
\end{frame}

\subsection*{5.35}
\begin{frame}
    \frametitle{题面}
    用正交变换法将二次型
    \begin{equation*}
        f\spar{x_1,x_2,x_3}=x_1^2+4x_2^2+4x_3^2-4x_1x_2+4x_1x_3-8x_2x_3
    \end{equation*}
    化为标准形，写出所作的正交变换。
\end{frame}

\begin{frame}[allowframebreaks]
    \frametitle{解答}
    记二次型\(f\)的矩阵为\(\mata=\begin{bmatrix}1&-2&2\\-2&4&-4\\2&-4&4\end{bmatrix}\)。则\(\mata\)的特征多项式为
    \begin{equation*}
        \abs{\lambda\mati-\mata}=\begin{vmatrix}\lambda-1&2&-2\\2&\lambda-4&4\\-2&4&\lambda-4\end{vmatrix}=\lambda^2\spar{\lambda-9}
    \end{equation*}
    所以\(\mata\)的特征值为\(\lambda_1=\lambda_2=0\)，\(\lambda_3=9\)。

    将上述特征值代入特征方程组\(\spar{\lambda\mati-\mata}\vecal=0\)，得到\(\lambda_1,\lambda_2\)对应的线性无关的特征向量为\(\vecal_1=\myvec{2,0,-1}\)和\(\vecal_2=\myvec{0,1,1}\)，\(\lambda_3\)对应的线性无关的特征向量为\(\vecal_3=\myvec{1,-2,2}\)。

    对向量组\(\vecal_1,\vecal_2\)和向量组\(\vecal_3\)分别做Schmidt正交化，得到正交变换的表示矩阵
    \begin{equation*}
        \matp=
        \begin{bmatrix}
            \frac{2\sqrt{5}}{5} & \frac{2\sqrt{39}}{39} & \frac{1}{3}  \\
            0                   & \frac{5\sqrt{39}}{39} & -\frac{2}{3} \\
            -\frac{\sqrt{5}}{5} & \frac{4\sqrt{39}}{39} & \frac{2}{3}
        \end{bmatrix}
    \end{equation*}
    即二次型\(f\)在正交变换\(\vecy=\matp^\top\vecx\)下化为标准形\(f\spar{x_1,x_2,x_3}=9y_3^2\)。
\end{frame}

\subsection*{5.36}
\begin{frame}
    \frametitle{题面}
    二次型
    \begin{equation*}
        f\spar{x_1,x_2,x_3}=\vecx^\top\mata\vecx=ax_1^2+2x_2^2-2x_3^2+2bx_1x_3\spar{b>0}
    \end{equation*}
    其中二次型\(f\)的矩阵\(\mata\)的特征值之和为\(1\)，特征值之积为\(-12\)，求
    \begin{enumerate}
        \item \(a,b\)的值；
        \item 用正交变换将\(f\)化为标准形，写出所作的正交变换及对应的正交矩阵。
    \end{enumerate}
\end{frame}

\begin{frame}[allowframebreaks]
    \frametitle{解答}
    \begin{enumerate}
        \item {
              由题意，\(\mata=\begin{bmatrix}a&0&b\\0&2&0\\b&0&-2\end{bmatrix}\)。由定理5.1.2，有
              \begin{equation*}
                  \begin{cases}
                      \trace{\mata}=a=1 \\
                      \det{\mata}=-4a-2b^2=-12
                  \end{cases}
              \end{equation*}
              解得\(a=1,b=2\)。
              }
        \item {
              \(\mata\)的特征多项式为
              \begin{equation*}
                  \abs{\lambda\mati-\mata}=\begin{vmatrix}\lambda-1&0&-2\\0&\lambda-2&0\\-2&0&\lambda+2\end{vmatrix}=\spar{\lambda-2}^2\spar{\lambda+3}
              \end{equation*}
              所以\(\mata\)的特征值为\(\lambda_1=\lambda_2=2\)，\(\lambda_3=-3\)。

              将上述特征值代入特征方程组\(\spar{\lambda\mati-\mata}\vecal=0\)，得到\(\lambda_1,\lambda_2\)对应的线性无关的特征向量为\(\vecal_1=\myvec{2,0,1}\)和\(\vecal_2=\myvec{0,1,0}\)，\(\lambda_3\)对应的线性无关的特征向量为\(\vecal_3=\myvec{1,0,-2}\)。

              对向量组\(\vecal_1,\vecal_2\)和向量组\(\vecal_3\)分别做Schmidt正交化，得到正交变换的表示矩阵
              \begin{equation*}
                  \matp=
                  \begin{bmatrix}
                      \frac{2\sqrt{5}}{5} & 0 & \frac{\sqrt{5}}{5}   \\
                      0                   & 1 & 0                    \\
                      \frac{\sqrt{5}}{5}  & 0 & -\frac{2\sqrt{5}}{5}
                  \end{bmatrix}
              \end{equation*}
              即二次型\(f\)在正交变换\(\vecy=\matp^\top\vecx\)下化为标准形\(f\spar{x_1,x_2,x_3}=2y_1^2+2y_2^2-3y_3^2\)。
              }
    \end{enumerate}
\end{frame}

\subsection*{5.38}
\begin{frame}
    \frametitle{题面}
    已知二次型\(f\spar{x_1,x_2,x_3}=4x_1^2+4x_2^2+4x_3^2+4x_1x_2+ax_1x_3-4x_2x_3\)的秩为\(2\)，
    \begin{enumerate}
        \item 求\(a\)的值；
        \item 求正交变换\(\vecx=\matp\vecy\)把\(f\spar{x_1,x_2,x_3}\)化为标准形；
        \item 求\(f\spar{x_1,x_2,x_3}=0\)的解。
    \end{enumerate}
\end{frame}

\begin{frame}[allowframebreaks]
    \frametitle{解答}
    记二次型\(f\)的矩阵为\(\mata=\begin{bmatrix}4&2&a/2\\2&4&2\\a/2&2&4\end{bmatrix}\)。
    \begin{enumerate}
        \item {
              因为\(\rank{\mata}=2\neq3\)，因此有
              \begin{equation*}
                  \det{\mata}=32+4a-a^2=0
              \end{equation*}
              解得\(a=8\)或\(a=-4\)，将\(a\)代入验证\(\rank{\mata}=2\)，均满足条件。
              }
        \item {
              当\(a=8\)时，\(\mata\)的特征多项式为
              \begin{equation*}
                  \abs{\lambda\mati-\mata}=\begin{vmatrix}\lambda-4&-2&-4\\-2&\lambda-4&-2\\-4&-2&\lambda-4\end{vmatrix}=\lambda\spar{\lambda^2-12\lambda+24}
              \end{equation*}
              所以\(\mata\)的特征值为\(\lambda_1=0\)，\(\lambda_2=6+2\sqrt{3}\)，\(\lambda_3=6-2\sqrt{3}\)。

              将上述特征值代入特征方程组\(\spar{\lambda\mati-\mata}\vecal=0\)，得到\(\mata\)的线性无关的特征向量为\(\vecal_1=\myvec{-1,0,1}\)，\(\vecal_2=\myvec{1,-1+\sqrt{3},1}\)，\(\vecal_3=\myvec{1,-1-\sqrt{3},1}\)。

              对\(\vecal_1,\vecal_2,\vecal_3\)标准化，得到正交变换的表示矩阵
              \begin{equation*}
                  \matp_1=
                  \begin{bmatrix}
                      -0.707 & 0.325  & -0.628 \\
                      0      & -0.888 & -0.460 \\
                      0.707  & 0.325  & -0.628
                  \end{bmatrix}
              \end{equation*}
              该正交变换将二次型\(f\)化为标准形\(f\spar{x_1,x_2,x_3}=\spar{6+2\sqrt{3}}y_2^2+\spar{6-2\sqrt{3}}y_3^2\)。

              当\(a=-4\)时，\(\mata\)的特征多项式为
              \begin{equation*}
                  \abs{\lambda\mati-\mata}=\begin{vmatrix}\lambda-4&-2&2\\-2&\lambda-4&-2\\2&-2&\lambda-4\end{vmatrix}=\lambda\spar{\lambda-6}^2
              \end{equation*}
              所以\(\mata\)的特征值为\(\lambda_1=\lambda_2=6\)，\(\lambda_3=0\)。

              将上述特征值代入特征方程组\(\spar{\lambda\mati-\mata}\vecal=0\)，得到\(\lambda_1,\lambda_2\)对应的线性无关的特征向量为\(\vecal_1=\myvec{1,1,0}\)和\(\vecal_2=\myvec{-1,0,1}\)，\(\lambda_3\)对应的线性无关的特征向量为\(\vecal_3=\myvec{1,-1,1}\)。

              对向量组\(\vecal_1,\vecal_2\)和向量组\(\vecal_3\)分别做Schmidt正交化，得到正交变换的表示矩阵
              \begin{equation*}
                  \matp_2=
                  \begin{bmatrix}
                      \frac{\sqrt{2}}{2} & -\frac{\sqrt{6}}{6} & \frac{\sqrt{3}}{3}  \\
                      \frac{\sqrt{2}}{2} & \frac{\sqrt{6}}{6}  & \frac{-\sqrt{3}}{3} \\
                      0                  & \frac{2\sqrt{6}}{6} & \frac{\sqrt{3}}{3}
                  \end{bmatrix}
              \end{equation*}
              该正交变换将二次型\(f\)化为标准形\(f\spar{x_1,x_2,x_3}=6y_1^2+6y_2^2\)。
              }
        \item {
              当\(a=8\)时，因为正交变换\(\matp_1\)将二次型\(f\)化为标准形\(f\spar{x_1,x_2,x_3}=\spar{6+2\sqrt{3}}y_2^2+\spar{6-2\sqrt{3}}y_3^2\)，因此\(f\spar{x_1,x_2,x_3}=0\)的解空间的基为
              \begin{equation*}
                  \matp_1\begin{bmatrix}1&0&0\end{bmatrix}^\top=\begin{bmatrix}-\frac{\sqrt{2}}{2}&0&\frac{\sqrt{2}}{2}\end{bmatrix}^\top
              \end{equation*}

              当\(a=-4\)时，因为正交变换\(\matp_2\)将二次型\(f\)化为标准形\(f\spar{x_1,x_2,x_3}=6y_1^2+6y_2^2\)，因此\(f\spar{x_1,x_2,x_3}=0\)的解空间的基为
              \begin{equation*}
                  \matp_2\begin{bmatrix}0&0&1\end{bmatrix}^\top=\begin{bmatrix}\frac{\sqrt{3}}{3}&-\frac{\sqrt{3}}{3}&\frac{\sqrt{3}}{3}\end{bmatrix}^\top
              \end{equation*}
              }
    \end{enumerate}
\end{frame}

\subsection*{5.39}
\begin{frame}
    \frametitle{题面}
    判断二次型的正定性：
    \begin{enumerate}
        \item \(f\spar{x_1,x_2,x_3}=5x_1^2+x_2^2+5x_3^2+4x_1x_2-8x_1x_3-4x_2x_3\)；
        \item \(f\spar{x_1,x_2,x_3,x_4}=x_1^2+x_2^2+14x_3^2+7x_4^2+6x_1x_3+8x_1x_4-4x_2x_3+2x_2x_4+4x_3x_4\)。
    \end{enumerate}
\end{frame}

\begin{frame}[allowframebreaks]
    \frametitle{解答}
    \begin{enumerate}
        \item {
              记二次型\(f\)的矩阵为\(\mata=\begin{bmatrix}5&2&-4\\2&1&-2\\-4&-2&5\end{bmatrix}\)。则其顺序主子式为
              \begin{equation*}
                  \Delta_1=\begin{vmatrix}5\end{vmatrix}=5>0,\Delta_2=\begin{vmatrix}5&2\\2&1\end{vmatrix}=1>0,
              \end{equation*}
              \begin{equation*}
                  \Delta_3=\begin{vmatrix}5&2&-4\\2&1&-2\\-4&-2&5\end{vmatrix}=1>0
              \end{equation*}
              所以该二次型是正定的。
              }
        \item {
              记二次型\(f\)的矩阵为\(\mata=\begin{bmatrix}1&0&3&4\\0&1&-2&1\\3&-2&14&2\\4&1&2&7\end{bmatrix}\)。因为\(\mata\)的\(4\)阶顺序主子式为
              \begin{equation*}
                  \Delta_4=\begin{vmatrix}1&0&3&4\\0&1&-2&1\\3&-2&14&2\\4&1&2&7\end{vmatrix}=-74<0
              \end{equation*}
              所以该二次型不是正定的。
              }
    \end{enumerate}
\end{frame}

\subsection*{5.40}
\begin{frame}
    \frametitle{题面}
    判断二次型\(f=\sum_{i=1}^nx_i^2+\sum_{1\le i<j\le n}x_ix_j\)的正定性。
\end{frame}

\begin{frame}
    \frametitle{解答}
    记二次型\(f_n\)的矩阵为\(\mata_n\in\rea^{n\times n}\)，则有
    \begin{gather*}
        A_n =
        \begin{bmatrix}
            1           & \frac{1}{2} & \cdots & \frac{1}{2} \\
            \frac{1}{2} & 1           & \cdots & \frac{1}{2} \\
            \vdots      & \vdots      &        & \vdots      \\
            \frac{1}{2} & \frac{1}{2} & \cdots & 1
        \end{bmatrix}\xrightarrow[i=2,3,\dots,n]{C_1+C_i}
        \begin{bmatrix}
            \frac{n+1}{2} & \frac{1}{2} & \cdots & \frac{1}{2} \\
            \frac{n+1}{2} & 1           & \cdots & \frac{1}{2} \\
            \vdots        & \vdots      &        & \vdots      \\
            \frac{n+1}{2} & \frac{1}{2} & \cdots & 1
        \end{bmatrix} \\
        \xrightarrow[i=2,3,\dots,n]{R_i-R_1}
        \begin{bmatrix}
            \frac{n+1}{2} & \frac{1}{2} & \cdots & \frac{1}{2} \\
            0             & \frac{1}{2} & \cdots & \frac{1}{2} \\
            \vdots        & \vdots      &        & \vdots      \\
            0             & 0           & \cdots & \frac{1}{2}
        \end{bmatrix}
    \end{gather*}
    \pause
    所以\(\det{\mata_n}=\spar{n+1}2^{-n}>0\)。又因为\(\mata_n\)的顺序主子式为\(\Delta_i=\det{\mata_i}\)，由定理5.4.2知\(\mata_n\)正定，即二次型\(f_n\)正定。
\end{frame}

\subsection*{5.41}
\begin{frame}
    \frametitle{题面}
    设\(\mata\)为\(3\)阶实对称矩阵，且满足\(\mata^2+2\mata=\mato\)，\(\rank{\mata}=2\)，
    \begin{enumerate}
        \item 求\(\mata\)的特征值；
        \item 当\(k\)为何值时，矩阵\(\mata+k\mati\)为正定矩阵？
    \end{enumerate}
\end{frame}

\begin{frame}
    \frametitle{解答}
    \begin{enumerate}
        \item {
              设\(\mata\)特征值为\(\lambda\)对应的特征向量为\(\vecal\neq\veczero\)，则有
              \begin{align*}
                  \spar{\lambda^2+2\lambda}\vecal & =\spar{\mata^2+2\mata}\vecal \\
                                                  & =\mato\vecal=\veczero
              \end{align*}
              所以\(\lambda^2+2\lambda=0\)，解得\(\lambda=0\)或\(\lambda=-2\)。
              }\pause
        \item {
              因为\(\rank{\mata}=2\)且\(\mata\)为实对称矩阵，因此存在正交矩阵\(\matp\)，满足\(\mata+k\mati=\matp\matlam\matp^\top\)，其中\(\matlam=\diag{\mat{k-2,k-2,k}}\)。由正定矩阵定义可知当\(k>2\)时，\(\mata+k\mati\)正定。
              }
    \end{enumerate}
\end{frame}

\subsection*{5.42}
\begin{frame}
    \frametitle{题面}
    设\(\mata\)是\(n\)阶实对称矩阵，证明\(\rank{\mata}=n\)的充要条件是存在一个\(n\)阶实矩阵\(\matb\)，使得\(\mata\matb+\matb^\top\mata\)正定。
\end{frame}

\begin{frame}
    \frametitle{解答}
    首先证明必要性。

    因为\(\rank{\mata}=n\)且\(\mata\)为实对称矩阵，可知\(\mata\)的特征值为非零实数\(\enums{\lambda}{n}\)。
    令\(\matb=\mata\)，则\(\mata\matb+\matb^\top\mata=2\mata^2\)特征值为\(2\lambda_1^2,2\lambda_2^2\dots,2\lambda_n^2\)。因此\(2\mata^2\)正定，结论成立。
    \pause

    随后用反证法证明充分性。

    因为\(\rank{\mata}<n\)，所以存在\(\vecx\in\rea^n\neq\veczero\)满足\(\mata\vecx=\veczero\)。因此对\(\forall\matb\in\rea^{n\times n}\)，有
    \begin{align*}
        \vecx^\top\spar{\mata\matb+\matb^\top\mata}\vecx & =\vecx^\top\spar{\mata\matb}\vecx+\vecx^\top\spar{\matb^\top\mata}\vecx \\
                                                         & =\spar{\mata\vecx}^\top\matb\vecx+\spar{\matb\vecx}^\top\mata\vecx=0
    \end{align*}
    这和题设矛盾。
\end{frame}

\subsection*{5.43}
\begin{frame}
    \frametitle{题面}
    若\(\mata\)是\(n\)阶实矩阵，证明：\(\trace{\mata\mata^\top}\ge0\)，且等号成立的充要条件是\(\mata=\mato\)。
\end{frame}

\begin{frame}
    \frametitle{解答}
    因为\(\mata\mata^\top\)为实对称矩阵，因此它可被对角化为\(\matp\matlam\matp^\top\)，其中\(\matp\)为正交矩阵。

    设\(\vece_i\in\rea^n\)表示第\(i\)个分量为\(1\)，其余分量为\(0\)的单位向量。因此得到
    \begin{align*}
        \vece_i^\top\matlam\vece_i & =\vece_i^\top\matp^\top\mata\mata^\top\matp\vece_i \\
                                   & =\Abs{\mata^\top\matp\vece_i}^2\ge0
    \end{align*}
    因此有\(\trace{\mata\mata^\top}=\sum_{i=1}^n\vece_i^\top\matlam\vece_i\ge0\)。
    \pause

    注意到当\(\vece_i^\top\matlam\vece_i=0\)对\(i=1,2,\dots,n\)均成立时等号成立，即\(\matlam=\mato\)。此时对\(\forall\vecx\in\rea^n\)，有
    \begin{equation*}
        \Abs{\mata^\top\vecx}^2=\vecx^\top\mata\mata^\top\vecx=\vecx^\top\matp\matlam\matp^\top\vecx=0
    \end{equation*}
    因此\(\mata^\top\vecx=\veczero\)对\(\forall\vecx\in\rea^n\)均成立，即\(\mata=\mata^\top=\mato\)。
\end{frame}

\subsection*{5.44}
\begin{frame}
    \frametitle{题面}
    设\(\mata\)是\(n\)阶正定矩阵，证明：对任意\(n\)阶矩阵\(\matb\)，\(\rank{\matb^\top\mata\matb}=\rank{\matb}\)。
\end{frame}

\begin{frame}
    \frametitle{解答}
    设\(\mata\)的特征值为\(\lambda_1>\lambda_2>\cdots>\lambda_n\)。因为\(\mata\)是正定矩阵，因此\(\mata\)可被对角化为\(\matp\matlam\matlam^\top\matp^\top\)，其中\(\matp\)为正交矩阵，\(\matlam=\diag{\mat{\sqrt{\lambda_1},\sqrt{\lambda_2},\dots,\sqrt{\lambda_n}}}\)。因此有
    \begin{equation*}
        \matb^\top\mata\matb=\matb^\top\matp\matlam\matlam^\top\matp^\top\matb=\spar{\matlam^\top\matp^\top\matb}^\top\spar{\matlam^\top\matp^\top\matb}
    \end{equation*}
    \pause
    所以有
    \begin{align*}
        \rank{\matb^\top\mata\matb} & =\rank{\spar{\matlam^\top\matp^\top\matb}^\top\spar{\matlam^\top\matp^\top\matb}} \\
                                    & =\rank{\matlam^\top\matp^\top\matb}=\rank{\matb}
    \end{align*}
\end{frame}

\end{document}